% Generated draft: requires mathematical and source review.
\documentclass[11pt]{article}
\usepackage[margin=25mm]{geometry}
\usepackage{fontspec}
\setmainfont{Latin Modern Roman}
\usepackage{amsmath,amssymb,mathtools,unicode-math}
\setmathfont{Latin Modern Math}
\usepackage{enumitem,needspace,xurl,hyperref}
\setlength{\parindent}{0pt}
\setlength{\parskip}{5pt}
\setlength{\emergencystretch}{3em}
\newcommand{\N}{\mathbb{N}}
\newcommand{\Z}{\mathbb{Z}}
\newcommand{\Q}{\mathbb{Q}}
\newcommand{\R}{\mathbb{R}}
\newcommand{\C}{\mathbb{C}}
\newcommand{\F}{\mathbb{F}}
\newcommand{\A}{\mathbb{A}}
\newcommand{\PP}{\mathbb{P}}
\newcommand{\E}{\mathbb{E}}
\newcommand{\eps}{\varepsilon}
\newcommand{\dd}{\,\mathrm{d}}
\begin{document}
\section*{UnsolvedMath ID: 800780}
\textbf{Source problem code:} catalogue entry
\subsection{Definitions and mathematical statement}
Fix a complex dimension \(n\geq 1\). A holomorphic local dynamical system at a point \(p\) of an \(n\)-dimensional complex manifold \(M,p)\) is understood here as a germ of a holomorphic map \(f:(M,p)\to(M,p)\) satisfying \(f(p)=p\). Equality of germs means equality on some neighborhood of the distinguished point. Write \(\mathcal H_n\) for the collection of such germs at \(0\in\C^n\). The supplied excerpt does not explicitly define local dynamical systems or impose invertibility, derivative conditions, or other restrictions; none of these additional restrictions is imposed here.

For a regularity category \(r\), a local conjugacy between \(f:(M,p)\to(M,p)\) and \(g\in\mathcal H_n\) is an invertible germ \(h:(M,p)\to(\C^n,0)\) satisfying
\[
 h\circ f=g\circ h.
\]
In the holomorphic category, \(h\) is biholomorphic. In the smooth category, \(h\) is a \(C^\infty\) diffeomorphism of the underlying real manifolds. In the \(C^k\) category, for a specified finite integer \(k\geq1\), \(h\) and \(h^{-1}\) are \(C^k\) maps of the underlying real manifolds. In the topological category, \(h\) is a homeomorphism. These are the usual interpretations of the listed categories; the source does not specify the permitted values of \(k\).

\textbf{Normal-form classification problem (Abate, Problem Q3).} For each chosen regularity category \(r\), find a possibly small family \(\mathcal F_{n,r}\subseteq\mathcal H_n\) such that every holomorphic local dynamical system \(f\) at every point \(p\) of every \(n\)-dimensional complex manifold admits a local conjugacy of category \(r\) to some \(g\in\mathcal F_{n,r}\). Such a \(g\) is called a normal form of \(f\) in that category. The families may depend on the regularity category.

The source asks for a \emph{possibly unique} normal form. Thus uniqueness is an optional strengthening, expressed by requiring that, for every \(f\),
\[
 g_1,g_2\in\mathcal F_{n,r},\quad
 f\sim_r g_1,\quad f\sim_r g_2
 \quad\Longrightarrow\quad g_1=g_2,
\]
where \(\sim_r\) denotes local conjugacy of category \(r\). This is uniqueness of the representative germ, not uniqueness of the conjugating map.

The excerpt supplies no mathematical criterion for how small, explicit, or practical the family must be. Consequently, it specifies a classification program without a quantitative or constructive success criterion beyond the covering property and optional uniqueness above. It cautions that holomorphic classification is often too complicated to be practical and that a normal-form family might be uncountable. Its closing suggestion to seek invariants instead of normal forms does not specify an additional invariant-classification problem.
\subsection{Short English statement}
Abate's Problem Q3 asks for normal forms representing all holomorphic map germs fixing a point in a chosen complex dimension, under holomorphic, smooth, \(C^k\), or topological local conjugacy. Unique representatives are optional. The source leaves the desired size and practical form of the classification unspecified and warns that holomorphic normal-form families may be uncountable.
\subsection{Sources}
\begin{itemize}
\item Local UnsolvedMath registry, permanent ID \(800780\), source problem code \texttt{catalogue entry}, title: \emph{Local Discrete Holomorphic Dynamics --- Open Problem Q3}. The supplied record attributes the problem to Abate, with proposed year \(2011\), and labels it \emph{open}; that label is supplied registry metadata, not an independently verified current mathematical assessment.
\item Supplied source list: Abate, \emph{Open Problems in Local Discrete Holomorphic Dynamics} (2011), Problem Q3. Supplied source URL: \url{https://arxiv.org/abs/1107.4876}. The record reports access on \texttt{2026-10-16} and extraction by \texttt{arxiv\_latex}; these are supplied provenance statements, not a claim of independent source access or verification.
\item The supplied registry record was created and updated at \texttt{2026-10-16T14:00:00Z}, with upstream revision \texttt{d371ee917ac9647fbb06cd6fc529deec6f3dda82}.
\end{itemize}
\end{document}
